\section*{الفصل \ref{ch:b3:origin-of-elements} --- من أين تأتي العناصر}
% Values from the Book 5 scratch calculator calc/ch30.py.

\begin{solution}{exo:b3:origin-of-elements:1}
(أ)~$\mathrm{Si/Fe} = 10^{7.51 - 7.46} = 10^{0.05} = 1.12$
و$\mathrm{O/Fe} = 10^{8.69 - 7.46} = 10^{1.23} = 17.0$: نحو ذرّةِ سيليكون
واحدة وسبع عشرة ذرّة أكسجين لكلّ ذرّة حديد. (ب)~$\mathrm{Fe/Au} =
10^{7.46 - 0.91} = 10^{6.55} = \num{3.55e6}$ و$\mathrm{Pb/Au} =
10^{1.95 - 0.91} = 10^{1.04} = 11.0$: بعد ذروة الحديد تنخفض الوفرات
بملايين المرّات، لكنّ الرصاص، على ذروته الخاصّة، أكثر وفرةً
من الذهب بأحد عشر ضعفًا.
\end{solution}

\begin{solution}{exo:b3:origin-of-elements:2}
(أ)~$10^{7.46 - 12} = \num{2.884e-5}$، أي \num{28.8} ذرّة حديد لكلّ
مليون ذرّة هيدروجين. (ب)~بمائة مرّة أقلّ: \num{0.288} لكلّ
مليون، ذرّة حديد واحدة لكلّ $1/\num{2.884e-7} = \num{3.47e6}$ ذرّة
هيدروجين. (ج)~$[\mathrm{Fe/H}] = \log_{10}(\num{e-7}/\num{2.884e-5}) =
-2.46$.
\end{solution}

\begin{solution}{exo:b3:origin-of-elements:3}
$^{124}\mathrm{Xe}$ و$^{126}\mathrm{Xe}$ يقعان على الجانب الغنيّ بالبروتونات،
بعيدًا عن مسار s، وهما محميّان بالنظيرين المستقرّين $^{124}\mathrm{Te}$
و$^{126}\mathrm{Te}$: \omterm{def:b3:origin-of-elements:p-process}{النوى p}. $^{128}\mathrm{Xe}$ و$^{130}\mathrm{Xe}$
يقعان على مسار s لكنّهما محميّان من \omterm{def:b3:origin-of-elements:s-process}{العملية r} بالنظيرين المستقرّين
$^{128}\mathrm{Te}$ و$^{130}\mathrm{Te}$: s فقط. $^{129}\mathrm{Xe}$،
$^{131}\mathrm{Xe}$ و$^{132}\mathrm{Xe}$ يقعان على مسار s ولا يملكان
نظيرًا مستقرًّا متساويًا في الكتلة وذا $Z$ أدنى (لليود 127 فقط): تُنتَجان بالعمليّتين معًا.
$^{134}\mathrm{Xe}$ و$^{136}\mathrm{Xe}$ يقعان بعد
$^{133}\mathrm{Xe}$، حيث يخرج مسار s من الزينون، وهما غير
محميّين: r فقط.
\end{solution}

\begin{solution}{exo:b3:origin-of-elements:4}
(أ)~$t = d/c = 40\times\qty{3.086e22}{m}/\qty{2.998e8}{m.s^{-1}} =
\qty{4.12e15}{s}$، أي نحو \num{130} مليون سنة. (ب)~$\Delta v/c \le
1.7/\num{4.12e15} = \num{4.1e-16}$. هذا حدّ أعلى، لأنّه قد يعود جزء من التأخير أو
كلّه إلى انبعاث أشعّة غاما بعد
الاندماج.
\end{solution}

\begin{solution}{exo:b3:origin-of-elements:5}
(أ)~$295\times0.1067 = \qty{31.48}{mb}$، $832\times0.0373 =
\qty{31.03}{mb}$، $155\times0.1984 = \qty{30.75}{mb}$. (ب)~انطلاقًا من
$^{122}\mathrm{Te}$: $N_{123} = 31.48/832 = \num{0.03783}$، \qty{1.4}{\percent}
أعلى من القيمة المقيسة، و$N_{124} = 31.48/155 = \num{0.2031}$،
\qty{2.4}{\percent} أعلى. الوفرات التي تختلف بعامل خمسة تُتنبَّأ
بدقّة في حدود بضعة في المئة من مقاطع التفاعل وحدها.
(ج)~$^{128}\mathrm{Te}$ يُنتَج \omterm{def:b3:origin-of-elements:s-process}{بالعملية r} فقط: يخرج دفق s من التيلوريوم عند
$^{127}\mathrm{Te}$، الذي يتفكّك، فلا يمرّ دفق s ثابت عبر
$^{128}\mathrm{Te}$، وقد تحدّدت وفرته \omterm{def:b3:origin-of-elements:s-process}{بالعملية r}.
\end{solution}

\begin{solution}{exo:b3:origin-of-elements:6}
(أ)~$\tau_0\langle\sigma\rangle = 0.306\times155 = 47.4$، إذن
$f = 1/(1 + 1/47.4) = \num{0.979}$ من أجل $^{124}\mathrm{Te}$؛ وأمّا من أجل
$^{138}\mathrm{Ba}$، $0.306\times4.0 = 1.224$ و$f = 1/(1 + 0.817) =
\num{0.550}$: تنقل النواة السحرية \qty{55}{\percent} فقط من
الدفق. (ب)~العوامل الستّة هي \num{0.959}، \num{0.550}، \num{0.908}،
\num{0.771}، \num{0.972} و\num{0.915}؛ وحاصل ضربها \num{0.328}:
$\sigma N$ ينخفض بعامل \num{3.05}. (ج)~الانخفاض المقيس هو
$27.6/7.72 = 3.57$. يعطي النموذج حجم القفزة بدقّة في حدود
\qty{15}{\percent}، ويأتي معظمها من $^{138}\mathrm{Ba}$
و$^{140}\mathrm{Ce}$، وهما النواتان ذواتا $N = 82$.
\end{solution}

\begin{solution}{exo:b3:origin-of-elements:7}
(أ)~$kT = 8\times10^3\times\qty{1.602e-19}{J} = \qty{1.282e-15}{J}$،
و$v_T = \sqrt{2\times\num{1.282e-15}/\num{1.6749e-27}} =
\qty{1.24e6}{m.s^{-1}}$. (ب)~بأخذ $\langle\sigma\rangle =
\qty{155e-31}{m^2}$، $\lambda = n_\mathrm{n}\langle\sigma\rangle v_T$
يساوي \qty{1.92e-11}{s^{-1}} من أجل $n_\mathrm{n} = \qty{e12}{m^{-3}}$
($10^6$ لكلّ سنتيمتر مكعّب): زمنٌ متوسّط قدره \qty{5.2e10}{s}، أي نحو
\num{1650} سنة؛ ثم \num{16.5} سنة عند $10^8$ و\num{0.165}~سنة،
أي شهرين، عند \qty{e10}{cm^{-3}}. (ج)~يتفكّك $^{127}\mathrm{Te}$ في
ساعات، أسرع بكثير من أيّ من أزمنة الأسر هذه: يتفكّك دائمًا، ولا
ينقسم المسار عنده. يعيش $^{134}\mathrm{Cs}$ نحو سنتين،
بين أزمنة الأسر عند $10^8$ و\qty{e10}{cm^{-3}}: وعند
الكثافات العالية لومضات الهيليوم يُؤسَر جزء كبير منه
قبل أن يتفكّك. وهو نقطة تفرّع، والوفرات على جانبيه
تقيس كثافة النيوترونات.
\end{solution}

\begin{solution}{exo:b3:origin-of-elements:8}
(أ)~$10^{6.30 - 3.14} = 1445$ مرّة أندر من الكالسيوم،
و$10^{4.97 - 3.14} = 68$ مرّة أندر من التيتانيوم. (ب)~$10^{8.46 - 1.38}
= \num{1.2e7}$. (ج)~يُظهر السكانديوم تأثير الفردي--الزوجي: النوى ذات
عدد بروتونات زوجيّ أشدّ ارتباطًا، وتُبنى
بأسر أنوية الهيليوم في مراحل الاحتراق في النجوم الضخمة، التي
تفضّلها. ويُظهر الليثيوم والبريليوم والبورون قاعًا في الوفرة:
هذه النوى الهشّة تُدمَّر عند درجات حرارة باطن النجوم، ويتحوّل
احتراق الهيليوم مباشرةً إلى الكربون،
دون أن يُنتجها.
\end{solution}

\begin{solution}{exo:b3:origin-of-elements:9}
(أ)~لتكن $\alpha$ \omterm{def:b3:origin-of-elements:alpha}{عناصر ألفا} و$\mathrm{Fe}_\mathrm{cc}$،
$\mathrm{Fe}_\mathrm{Ia}$ الحديد الناتج عن نوعَي المستعرات الأعظم. على
الهضبة، $\alpha/\mathrm{Fe}_\mathrm{cc} = 10^{P}$ مرّة النسبة الشمسية
لها. عند $[\mathrm{Fe/H}] = 0$، فإنّ $\alpha/(\mathrm{Fe}_\mathrm{cc} +
\mathrm{Fe}_\mathrm{Ia})$ شمسيةٌ. وبالقسمة،
$\mathrm{Fe}_\mathrm{cc}/(\mathrm{Fe}_\mathrm{cc} + \mathrm{Fe}_\mathrm{Ia})
= 10^{-P}$، وحصّة النوع \mbox{\textdir TLT Ia} هي $1 - 10^{-P} = 1 - 10^{-0.4} =
\num{0.602}$. (ب)~$10^{-P} = 1/3$ يعطي $P = \log_{10}3 = 0.48$.
(ج)~تبدأ المستعرات الأعظم من النوع \mbox{\textdir TLT Ia} بعد التأخير نفسه في المجرّتين.
فحين بدأت في المجرّة القزمة، لم يبلغ غازُها سوى
$[\mathrm{Fe/H}] = -2$، مقابل $-1$ في درب التبانة: كانت قد حوّلت غازها
إلى نجوم بوتيرة أبطأ بعشر مرّات.
\end{solution}

\begin{solution}{exo:b3:origin-of-elements:10}
(أ)~بالنسبة إلى الزركونيوم،
\[
\frac{(X/\mathrm{Mg})_\mathrm{Earth}}{(X/\mathrm{Mg})_\mathrm{CI}} =
\frac{10.5/228000}{3.96/95870} = 1.11 ;
\]
وبالمثل $(2670/228000)/(5010/95870) = 0.224$ للصوديوم،
\num{0.190} للبوتاسيوم و\num{0.118} للروبيديوم. (ب)~الزركونيوم، الأشدّ امتناعًا عن التطاير، يقع
قريبًا من نسبته في الكوندريت. وتُستنفد الفلزات القلوية الثلاثة المتوسطة التطاير،
التي تتكثّف عند \qtyrange{800}{1006}{K}، بعوامل من خمسة إلى تسعة،
والروبيديوم، الأكثر تطايرًا، الأكثر استنفادًا. وتتبع
الأرض السيليكاتية \omterm{def:b3:origin-of-elements:volatility}{اتجاه التطاير}. (ج)~المغنيسيوم عنصر رئيس ممتنع عن التطاير،
يتكثّف مع معظم الصخر، وهو
محبّ للصخر: بقي في الوشاح حين تكوّن اللبّ. ولو اتُّخذ مرجعٌ محبّ للحديد
مثل الحديد أو النيكل، لخلط الفقد إلى اللبّ مع
الفقد بالتطاير.
\end{solution}

\begin{solution}{exo:b3:origin-of-elements:11}
(أ)~عندما تُحبس كتلة $\mathrm{d}s$ في النجوم، تزيل الرياح
$\lambda\,\mathrm{d}s$ من الغاز ذي \omterm{def:b3:origin-of-elements:metallicity}{المعدنية} $Z$:
$\mathrm{d}M_\mathrm{g} = -(1 + \lambda)\,\mathrm{d}s$
و$\mathrm{d}(ZM_\mathrm{g}) = y\,\mathrm{d}s - Z\,\mathrm{d}s -
\lambda Z\,\mathrm{d}s$. بنشر الحدّ الثاني وطرح
$Z\,\mathrm{d}M_\mathrm{g}$ فيبقى $M_\mathrm{g}\,\mathrm{d}Z =
y\,\mathrm{d}s = -\frac{y}{1 + \lambda}\,\mathrm{d}M_\mathrm{g}$، وهذا
يُكامَل إلى النتيجة. (ب)~$\ln(1/0.29) = 1.238$: $Z = 0.0155$ عند
$\lambda = 0$، ونصفها، $0.0077$، عند $\lambda = 1$. (ج)~تُخفّض الرياح
المردود الفعّال بعامل $1 + \lambda$. وفي مجرّة صغيرة الكتلة،
تدفع المستعرات الأعظم الغاز إلى الخارج بسهولة أكبر، و$\lambda$ كبيرة،
فتُفقد الفلزات مع الغاز.
\end{solution}

\begin{solution}{exo:b3:origin-of-elements:12}
(أ)~$N_\mathrm{Li}/N_\mathrm{H} = 10^{0.96 - 12} = \num{9.1e-12}$ في
الشمس، و$10^{3.25 - 12} = \num{1.78e-9}$ في كوندريتات CI. وتحتفظ الشمس
بنسبة $10^{-2.29} = \num{0.0051}$ من ليثيومها: فقد دمّرت
\qty{99.5}{\percent} منه. (ب)~$\num{1.78e-9}/\num{1.45e-10} = 12.3$
و$\num{1.78e-9}/\num{4.72e-10} = 3.8$. (ج)~تحتوي المجموعة الشمسية على قرابة
أربعة أضعاف ما صنعه الانفجار الكبير من الليثيوم: معظم ليثيومها
صُنع لاحقًا بعملية أخرى. وأقدم النجوم تحتوي على ما يقلّ بثلاث مرّات
عمّا صنعه الانفجار الكبير: إمّا أنّها دمّرت جزءًا منه، كما فعلت
الشمس، أو أنّ التنبؤ يغفل شيئًا؛ والمقارنتان
وحدهما لا تبيّنان أيّهما.
\end{solution}

\begin{solution}{pb:b3:origin-of-elements:1}

\textbf{الجزء الأول.}

\textbf{1.} $10^{12 - 7.46} = 10^{4.54} = \num{34700}$.

\textbf{2.} $N_\mathrm{He}/N_\mathrm{H} = 10^{10.914 - 12} = \num{0.0820}$، إذن
$4N_\mathrm{He}/N_\mathrm{H} = \num{0.328}$، مقابل $Y/X =
0.2423/0.7438 = \num{0.326}$. يتّفقان إلى حدود \qty{0.7}{\percent}؛ ويرجع
الفرق الصغير إلى أنّ ذرّة الهيليوم تزن أقل قليلًا من أربع
ذرّات هيدروجين.

\textbf{3.} $\log_{10}\bigl((1/350000)/\num{2.884e-5}\bigr) = -5.544 + 4.540 =
-1.00$.

\textbf{4.} $[\mathrm{O/H}] = -1.00 + 0.40 = -0.60$، إذن
$\log\varepsilon(\mathrm{O}) = 8.69 - 0.60 = 8.09$. يحتوي النجم على
$10^{8.09 - 6.46} = 42.7$ ذرّة أكسجين لكلّ ذرّة حديد، مقابل
$10^{1.23} = 17.0$ في الشمس: أي أكثر بمقدار $10^{0.4} = 2.5$ مرّة.

\textbf{5.} $Z \approx 0.0139\times10^{-0.60} = \num{0.0035}$، وهو ربع
قيمة الشمس.

\textbf{6.} للحديد عدد كبير جدًّا من الخطوط في طيف نجم بارد، بينما
للأكسجين عدد قليل منها وخافت: يُقاس الحديد بسهولة أكبر بكثير
وبدقّة أعلى، ويُتّخذ بديلًا عن الفلزات.

\textbf{الجزء الثاني.}

\textbf{7.} يقع كلٌّ منها خلف خطّ تساوي الضغط المستقرّ الذي له $Z$ أدنى:
$^{122}\mathrm{Sn}$، و$^{123}\mathrm{Sb}$ و$^{124}\mathrm{Sn}$ توقف
سلاسل التحلّل في \omterm{def:b3:origin-of-elements:s-process}{العملية r} قبل أن تبلغ التيلوريوم. وتقع النوى الثلاثة
على مسار s، وهو يمرّ بالتيلوريوم من $^{122}\mathrm{Te}$ إلى
$^{126}\mathrm{Te}$.

\textbf{8.} \num{31.48}، \num{31.03} و\qty{30.75}{mb}؛ المتوسط
\qty{31.09}{mb}.

\textbf{9.} متوسط الاثنين الآخرين هو $(31.48 + 30.75)/2 = \qty{31.11}{mb}$،
إذن $N_{123} = 31.11/832 = \num{0.03740}$، بفارق \qty{0.3}{\percent} فقط
فوق القيمة المقاسة \num{0.0373}.

\textbf{10.} $1.329\times44.4 = \qty{59.0}{mb}$، أي ما يقارب ضعف الهضبة.
$^{128}\mathrm{Te}$ تُنتَج \omterm{def:b3:origin-of-elements:s-process}{بالعملية r} وحدها: يغادر مسار s التيلوريوم عند
$^{127}\mathrm{Te}$، لذا صُنعت وفرتُها \omterm{def:b3:origin-of-elements:s-process}{بالعملية r}
ولا تدين بشيء لتدفّق \omterm{def:b3:origin-of-elements:s-process}{العملية s}.

\textbf{11.} $1 - 7.43/8.53 = \num{0.13}$: \qty{13}{\percent} من التدفّق مرّ
حول $^{148}\mathrm{Sm}$، بالأسر على النوى غير المستقرّة
$^{147}\mathrm{Nd}$، و$^{147}\mathrm{Pm}$ و$^{148}\mathrm{Pm}$ قبل
أن تتحلّل إلى السماريوم.

\textbf{12.} مع غلاف مغلق من 82 نيوترونًا، تربط هذه النوى نيوترونًا إضافيًّا
بضعف، ولها مقاطع أسر صغيرة جدًّا، إذ تبلغ \qty{4.0}{mb}
في حالة $^{138}\mathrm{Ba}$. يتراكم التدفّق عليها ولا يمضي منه إلّا جزء:
فهي عنق زجاجة.

\textbf{الجزء الثالث.}

\textbf{13.} $\mu_1 = 13.7/(13.7 + 33.4) = 13.7/47.1 = \num{0.291}$.

\textbf{14.} $\ln(1/0.291) = \num{1.235}$، إذن $y = 0.0154/1.235 = \num{0.0125}$،
\num{0.81} مرّة $Z_\odot$.

\textbf{15.} عند $\mu = 0.9$، يكون $\ln(1/0.9) = 0.1054$، إذن $Z = 0.0125\times0.1054 =
\num{0.00131}$. بما أنّ $0.00131/0.0154 = 0.0851$، فإنّ $[\mathrm{Fe/H}] =
-1.07$. وعند $\mu = 0.5$، يكون $Z = 0.0125\times0.693 = \num{0.00864}$
و$[\mathrm{Fe/H}] = -0.25$.

\textbf{16.} $Z = 0.00154$، إذن $\mu = \mathrm{e}^{-0.00154/0.01247} =
\mathrm{e}^{-0.1235} = \num{0.884}$.

\textbf{17.} $1 - 0.884 = \num{0.116}$: \qty{12}{\percent} من الكتلة.

\textbf{18.} $y = 0.0154/\ln 10 = \num{0.0067}$، أي نحو نصفها.

\textbf{الجزء الرابع.}

\textbf{19.} كتلة النجوم ذات \omterm{def:b3:origin-of-elements:metallicity}{المعدنية} دون $Z$ هي
$M(1 - \mathrm{e}^{-Z/y})$، إذن لكل وحدة من $Z$ تكون
$(M/y)\,\mathrm{e}^{-Z/y}$. بما أنّ $Z = Z_\odot10^{[\mathrm{Fe/H}]}$،
فإنّ $\mathrm{d}Z = Z\ln 10\,\mathrm{d}[\mathrm{Fe/H}]$، ولكل وحدة
من $[\mathrm{Fe/H}]$ يتناسب التوزيع مع
$Z\mathrm{e}^{-Z/y}$. مشتقّته،
$(1 - Z/y)\mathrm{e}^{-Z/y}$، تنعدم عند $Z = y$: تقع القمّة عند
$[\mathrm{Fe/H}] = \log_{10}0.810 = -0.09$.

\textbf{20.} $\num{0.116}/(1 - 0.291) = 0.1162/0.7091 = \num{0.164}$.

\textbf{21.} $Z/y = 1.235\times10^{-0.5} = 0.3905$, $1 - \mathrm{e}^{-0.3905} =
0.3233$، و$0.3233/0.7091 = \num{0.456}$.

\textbf{22.} $6670\times0.1638 = 1093$، مقابل 26: أي أكثر بأربعين مرّة من اللازم.

\textbf{23.} $26/902 = \qty{2.9}{\percent}$، وهي ما زالت أدنى بأكثر من خمس مرّات
من الصندوق. نجوم النموذج هي النجوم طويلة العمر، أي التي لا تتجاوز كتلتها
نحو كتلة شمسية واحدة، وقد بقيت جميعها منذ تكوّنها؛ أمّا من بين
نجوم F الأكثر كتلة، فقد فنيت أقدمها، وهي أيضًا
الأشدّ فقرًا بالمعادن، والعيّنة الكاملة تُقلّل من عددها.

\textbf{24.} $(1 - 0.1^{0.1})/(1 - 0.1) = 0.2057/0.9 = \num{0.23}$: أسوأ.
ولا يفيد رفع كسر الغاز كذلك: فمع اتجاه $\mu_1$ نحو 1
يتجه الكسر نحو $0.1$، وهو ما زال أكثر من ثلاث مرّات
من \qty{2.9}{\percent} المرصودة. لا يُنقذ أيّ كسر للغاز \omterm{thm:b3:origin-of-elements:closed-box}{نموذج الصندوق المغلق}.

\textbf{25.} يتنبّأ \omterm{thm:b3:origin-of-elements:closed-box}{نموذج الصندوق المغلق} بأنّ
\textbf{\qty[text-series-to-math=true]{16}{\percent}} من النجوم طويلة العمر
في الجوار الشمسي لها $[\mathrm{Fe/H}] < -1$، مقابل
\qty{3}{\percent} على الأكثر في الرصد. لم يكن القرص صندوقًا مغلقًا: فقد
نما بتدفّق بطيء من غاز فقير بالمعادن، بحيث قلّما
بقي غاز منتظرًا عند معدنية منخفضة لكي يكوّن النجوم.
\end{solution}
